\clearpage\thispagestyle{empty}
\pdfbookmark[0]{Quellen zum Pallas-Bericht vom Juli 1802}{pallas180207-sources}
\begin{center}\textbf{Quellen zum Pallas-Bericht vom Juli 1802}\end{center}
\noindent Diese KI-Übertragung in heutiges Deutsch folgt dem deutschen Auszug in Carl Friedrich Gauß’ \href{https://gdz.sub.uni-goettingen.de/id/PPN236007467}{\textit{Werken}, Band~VI} (Göttingen, 1874), S.~217--220. Franz Xaver von Zach berichtet über die von Heinrich Wilhelm Matthias Olbers entdeckte Pallas. Der Auszug enthält Rechnungen von Gauß und zitiert seinen Brief vom 26. Juni.
\par\smallskip\noindent Von Zachs \href{https://api.digitale-sammlungen.de/iiif/image/v2/bsb10538598_00077/full/1450,/0/default.jpg}{ursprünglicher Aufsatz}, Nr.~X, erschien in der \textit{Monatlichen Correspondenz zur Beförderung der Erd- und Himmels-Kunde} (Gotha, Becker), Band~VI, Juli 1802, S.~71--96. Der Auszug beginnt mit dem \href{https://api.digitale-sammlungen.de/iiif/image/v2/bsb10538598_00086/full/1450,/0/default.jpg}{Brief auf S.~80}, setzt mit den dritten Bahnelementen auf S.~83 wieder ein und endet mit der \href{https://api.digitale-sammlungen.de/iiif/image/v2/bsb10538598_00100/full/1450,/0/default.jpg}{Erörterung der Durchmesser auf S.~94--95}.
\par\smallskip\noindent Im \href{https://api.digitale-sammlungen.de/iiif/image/v2/bsb10538598_00090/full/1450,/0/default.jpg}{Zeitschriftendruck, S.~84}, wird für die mittlere Bewegung der Ceres auf die Elemente VII verwiesen. Die Unterschiede der Rektaszension am 5. und 7. April sind dort \(+0{,}7\) und \(0\). In den \textit{Werken} stehen IV sowie \(+0.\) und \(0.7\). Beide Übertragungen folgen den \textit{Werken}.
\par\smallskip\noindent Der im Brief angeführte Aufsatz von \href{https://api.digitale-sammlungen.de/iiif/image/v2/bsb10538597_00562/full/1450,/0/default.jpg}{Prof.\ Wurm} behandelt die Planetenmassen (Juni 1802, S.~546--570). Der eigenständige Aufsatz von \href{https://doi.org/10.1098/rstl.1802.0010}{William Herschel} heißt \textit{Observations on the Two Lately Discovered Celestial Bodies} (1802, S.~213--232).
\par\smallskip\noindent Übertragung und Satz wurden mit KI erstellt. Die Digitalisate der Zeitschrift stammen von der Bayerischen Staatsbibliothek.
\clearpage
